%% UML — Diagrama de casos de uso (plantilla)
%% Compilar: tectonic -X compile uml-casos-uso.tex
\documentclass[11pt,a4paper]{article}
\usepackage{../../docs/latex/legasoft}
\usepackage{../../docs/latex/legasoft-diagramas}

\lgtitulo{UML: diagrama de casos de uso}
\lgtituloCorto{UML — Casos de uso}
\lgcodigo{LGS-MOD-UML-1}
\lgversion{0.1}
\lgestado{Plantilla}
\lgfecha{\campo{fecha}}
\lgautor{\lgDirector\ (\lgDirectorRol)}

\begin{document}

\begin{center}
{\sffamily\bfseries\LARGE\color{lgPrimario} UML: diagrama de casos de uso\par}
\vspace{2pt}
{\sffamily\normalsize\color{lgGris} \campo{nombre del sistema} · Legasoft\par}
\end{center}
\vspace{6pt}
{\color{lgBorde}\hrule height 0.8pt}
\vspace{10pt}

\begin{porcompletar}[Cómo usar esta plantilla]
Cada caso de uso debe corresponder a uno o varios requisitos funcionales de
\texttt{LGS-REQ-001}. Un caso de uso expresa un objetivo del actor, no una
pantalla ni una operación técnica.
\end{porcompletar}

% =============================================================================
\section*{Diagrama}

\begin{figure}[H]
\centering
\ajustaancho{%
\begin{tikzpicture}[uml]

  % --- Casos de uso ---------------------------------------------------------
  \node[umlcaso] (cu1) {\campo{CU-01}\\ \campo{Registrar información}};
  \node[umlcaso, below=0.9cm of cu1] (cu2) {\campo{CU-02}\\ \campo{Consultar registros}};
  \node[umlcaso, below=0.9cm of cu2] (cu3) {\campo{CU-03}\\ \campo{Generar reporte}};
  \node[umlcaso, below=0.9cm of cu3] (cu4) {\campo{CU-04}\\ \campo{Administrar usuarios}};

  % Validar credenciales, incluido por otros casos
  \node[umlcaso, right=2.6cm of cu2, text width=2.6cm, align=center,
         minimum width=2.9cm]
    (cu5) {\campo{CU-05}\\ \campo{Autenticarse}};

  % Límite del sistema
  \begin{scope}[on background layer]
    \node[umlsistema, fit=(cu1)(cu2)(cu3)(cu4)(cu5),
          label={[font=\sffamily\small\bfseries, text=lgPrimario,
                  anchor=north]north:{\campo{Nombre del sistema}}}] {};
  \end{scope}

  % --- Actores --------------------------------------------------------------
  \umlactor{-4.2}{-1.35}{\campo{Usuario}\\\campo{operativo}}
  \coordinate (a1) at (-4.2,-1.35);

  \umlactor{-4.2}{-4.5}{\campo{Administrador}}
  \coordinate (a2) at (-4.2,-4.5);

  % --- Asociaciones ---------------------------------------------------------
  \draw[umlasoc] (a1) ++(0.35,0) -- (cu1.west);
  \draw[umlasoc] (a1) ++(0.35,0) -- (cu2.west);
  \draw[umlasoc] (a1) ++(0.35,0) -- (cu3.west);
  \draw[umlasoc] (a2) ++(0.35,0) -- (cu4.west);
  \draw[umlasoc] (a2) ++(0.35,0) -- (cu3.west);

  % --- Relación «include» ---------------------------------------------------
  \draw[umlext] (cu1.east) -- (cu5.north west)
    node[pos=0.5, above, sloped, font=\sffamily\scriptsize] {$\ll$include$\gg$};
  \draw[umlext] (cu4.east) -- (cu5.south west)
    node[pos=0.5, below, sloped, font=\sffamily\scriptsize] {$\ll$include$\gg$};

\end{tikzpicture}}
\caption{Casos de uso del sistema y los actores que los inician.}
\end{figure}

% =============================================================================
\Needspace*{20\baselineskip}
\section*{Especificación resumida de casos de uso}

\begin{center}
\begin{tabularx}{\linewidth}{@{}C{1.7cm} L{3.4cm} L{2.6cm} Y@{}}
\toprule
\sffamily\bfseries\color{lgPrimario}ID &
\sffamily\bfseries\color{lgPrimario}Caso de uso &
\sffamily\bfseries\color{lgPrimario}Actor principal &
\sffamily\bfseries\color{lgPrimario}Requisito asociado \\
\midrule
CU-01 & \campo{Registrar información} & \campo{Usuario operativo} & \campo{RF-001} \\
CU-02 & \campo{Consultar registros}   & \campo{Usuario operativo} & \campo{RF-002} \\
CU-03 & \campo{Generar reporte}       & \campo{Usuario / Admin}   & \campo{RF-003} \\
CU-04 & \campo{Administrar usuarios}  & \campo{Administrador}     & \campo{RF-004} \\
CU-05 & \campo{Autenticarse}          & \campo{Todos}             & \campo{RNF-002} \\
\bottomrule
\end{tabularx}
\captionof{table}{Casos de uso y su trazabilidad con los requisitos.}
\end{center}

\begin{nota}[Notación]
La línea continua representa una asociación entre actor y caso de uso. La línea
discontinua con $\ll$include$\gg$ indica que un caso de uso incorpora
obligatoriamente el comportamiento de otro. Para comportamiento opcional se usa
$\ll$extend$\gg$ en sentido inverso.
\end{nota}

\end{document}
